\documentclass{beamer}

\usepackage[spanish]{babel}
\decimalpoint

\usepackage{graphicx}
\usepackage{algorithm}
\usepackage{algpseudocode}
\usepackage{xcolor}
\usepackage{tikz}
\usepackage{booktabs}
\usepackage{bm}
\usepackage{amsmath}

\usetheme{Warsaw}
\usepackage{styles/warsaw-mod}
\usepackage{styles/warsaw-dark}

% Configuración de colores para el código
\usepackage{listings}
\lstset{
language=Python,
basicstyle=\ttfamily\scriptsize,
keywordstyle=\color{blue},
stringstyle=\color{green!50!black},
commentstyle=\color{gray},
showstringspaces=false,
frame=single,
extendedchars=true,
literate={á}{{'a}}1 {é}{{'e}}1 {í}{{'i}}1 {ó}{{'o}}1 {ú}{{'u}}1 {Á}{{'A}}1 {É}{{'E}}1 {Í}{{'I}}1 {Ó}{{'O}}1 {Ú}{{'U}}1 {ñ}{{~n}}1 {Ñ}{{~N}}1
}

\setbeamertemplate{navigation symbols}{}

% Cargar los metadatos de la presentación
\title{Algoritmos de Optimización de Funciones}
\subtitle{Práctica Calificada 2 - Métodos de primer y segundo orden}
\author[Gallo, Luna, Aguilar, Almendrades]{
    \begin{tabular}{r@{ -- }l}
        Adrian Marcelo Gallo Santillan & 20212022B \\
        Aldo Luna Bueno & 20170325B \\
        Víctor Manuel Aguilar Vitancio & 20212196K \\
        Jhoel Jhonny Almendrades Osorio & 20185513D
    \end{tabular}
}
\institute[UNI]{
    Universidad Nacional de Ingeniería\\
    Facultad de Ciencias - Ciencia de la Computación\\
    Matemática Computacional (CC3M2)
}
\date{\today}
% Descomentar y ajustar la ruta de la imagen cuando tengas el logo:
% \titlegraphic{\includegraphics[width=2cm]{logo_uni.png}}

\begin{document}

% Diapositiva de título
\begin{frame}
    \titlepage
\end{frame}

\begin{frame}{Contenido}
\tableofcontents
\end{frame}

% A. First-Order Methods
\section{First-Order Methods}
\subsection{Conjugate Gradient}

\begin{frame}{Conjugate Gradient: Explicación Gráfica}
    \begin{block}{Concepto principal}
        [Explica aquí la idea central de direcciones conjugadas y ortogonalidad respecto a la matriz hessiana].
    \end{block}
    \begin{itemize}
        \item {[Idea gráfica 1: Comportamiento en curvas de nivel cuadráticas frente al descenso estándar].}
        \item {[Idea gráfica 2: Ortogonalidad de las direcciones de búsqueda sucesivas].}
    \end{itemize}
\end{frame}

\begin{frame}{Conjugate Gradient: Pseudocódigo}
    \begin{algorithmic}[1]
        \State \textbf{Input:} {[Definir entradas: función \(f\), gradiente \(\nabla f\), punto inicial \(x^{(1)}\), tolerancias]}
        \State Inicializar dirección de descenso inicial \(d^{(1)}\)
        \While{[Condición de parada]}
            \State Calcular longitud de paso \(\alpha\)
            \State Actualizar punto: \(x \leftarrow x + \alpha d\)
            \State Calcular factor de conjugación \(\beta\)
            \State Actualizar nueva dirección: \(d \leftarrow -\nabla f + \beta d\)
        \EndWhile
        \State \textbf{Return} {[Punto óptimo]}
    \end{algorithmic}
\end{frame}

\begin{frame}{Conjugate Gradient: Ejemplo Paso a Paso}
    \begin{exampleblock}{Planteamiento del problema}
        [Define aquí una función de prueba cuadrática simple, ej: \(f(x_1, x_2) = x_1^2 + 2x_2^2\), y punto inicial].
    \end{exampleblock}
    \begin{itemize}
        \item \textbf{Paso 1:} {[Cálculo del gradiente inicial y primera dirección de búsqueda].}
        \item \textbf{Paso 2:} {[Cálculo del factor de actualización \(\beta\) y siguiente dirección].}
        \item \textbf{Paso 3:} {[Estado de convergencia y resultado].}
    \end{itemize}
\end{frame}

\begin{frame}[fragile]{Conjugate Gradient: Código}
\begin{lstlisting}[language=Python]
def conjugate_gradient(f, grad_f, x0, max_iter=50, tol=1e-5):
    # TODO: Inicializar variables y primera direccion d
    # while iter < max_iter:
    #     alpha = line_search(...)
    #     x = x + alpha * d
    #     beta = compute_beta(...)
    #     d = -grad + beta * d
    pass
\end{lstlisting}
\end{frame}

\begin{frame}{Conjugate Gradient: Análisis de Convergencia}
    \begin{alertblock}{Tasa de Convergencia}
        [Indicar tipo de convergencia en funciones cuadráticas vs funciones no lineales].
    \end{alertblock}
    \begin{itemize}
        \item {[Propiedades teóricas de convergencia en \(n\) dimensiones].}
        \item {[Ventajas y limitaciones frente a métodos de segundo orden puros].}
    \end{itemize}
\end{frame}
\subsection{Momentum}

% Paleta de curvas de nivel (misma que en Conjugate Gradient)
\colorlet{niv1}{AmarilloAviso!85!black}
\colorlet{niv2}{AmarilloAviso!50!VerdeNeon!75!black}
\colorlet{niv3}{VerdeNeon!65!black}
\colorlet{niv4}{VerdeNeon!40!CelesteBrillante!65!black}
\colorlet{niv5}{CelesteBrillante!65!black}
\colorlet{niv6}{CelesteBrillante!48!black}
\colorlet{niv7}{CelesteBrillante!34!black}
\providecommand{\reflibro}[1]{\par\vfill{\tiny\textcolor{GrisTexto!70}{\textit{Ref. libro:} #1}}}
% Curvas de nivel de f = 0.5 x1^2 + 5 x2^2 (ejes r y r/sqrt(10))
\newcommand{\momElipses}{%
    \foreach \r/\cl in {0.6/niv1, 1.3/niv2, 2.0/niv3, 2.8/niv4, 3.6/niv5, 4.5/niv6}
        \draw[\cl] (0,0) ellipse ({\r} and {\r/3.1623});%
}

% ---------------------------------------------------------------
% 1. Explicación gráfica: un mismo problema, iteración a iteración
%    Valle girado 30 grados, kappa=10, x(1)=(-3.6,-1.2), alpha=0.17, beta=0.5
% ---------------------------------------------------------------
\newcommand{\momValle}{%
    \foreach \r/\cl in {0.6/niv1, 1.2/niv2, 1.9/niv3, 2.6/niv4, 3.4/niv5, 4.3/niv6, 5.3/niv7}
        \draw[\cl, rotate around={30:(0,0)}] (0,0) ellipse ({\r} and {\r/3.1623});%
}

\begin{frame}{Momentum: Explicación Gráfica (1/3)}
    \begin{block}{Concepto principal}
        \footnotesize
        Una \textbf{bola que rueda}: la velocidad $v$ acumula los gradientes pasados.
        \[ v^{(k+1)}=\beta\,v^{(k)}-\alpha\,g^{(k)}, \qquad x^{(k+1)}=x^{(k)}+v^{(k+1)} \]
    \end{block}

    \begin{columns}[T]
        \begin{column}{0.56\textwidth}
            \centering
            \begin{tikzpicture}[x=1.3cm, y=1.3cm, font=\scriptsize]
                \begin{scope}
                    \clip (-3.9,-2.6) rectangle (-0.4,0.2);
                    \momValle
                    \draw[NaranjaBase, very thick, -stealth] (-3.6,-1.2) -- (-2.43,-1.99);
                \end{scope}
                \draw[gray!50] (-3.9,-2.6) rectangle (-0.4,0.2);
                \fill[GrisTexto] (-3.6,-1.2) circle (1.8pt) node[above, inner sep=2pt] {$x^{(1)}$};
                \fill[NaranjaBase] (-2.406,-2.004) circle (1.8pt) node[right, inner sep=3pt, text=GrisTexto] {$x^{(2)}$};
                \node[NaranjaBase, anchor=north east] at (-3.0,-1.65) {$v^{(2)}$};
                \draw[-stealth, GrisTexto!70] (-1.2,-0.55) -- (-0.75,-0.29);
                \node[GrisTexto!80, anchor=north] at (-1.0,-0.55) {hacia $x^*$};
            \end{tikzpicture}
        \end{column}
        \begin{column}{0.42\textwidth}
            \begin{block}{Iteración 1}
                \footnotesize
                $v^{(1)}=0$ $\Rightarrow$ $v^{(2)}=-\alpha\,g^{(1)}$
            \end{block}
            \small
            \begin{itemize}
                \item Sin velocidad previa: \textbf{igual que GD}
                \item El gradiente apunta casi \textbf{a través} del valle, no hacia $x^*$
            \end{itemize}
        \end{column}
    \end{columns}
    \reflibro{Ec.~(5.20)--(5.21).}
\end{frame}

\begin{frame}{Momentum: Explicación Gráfica (2/3)}
    \begin{columns}[T]
        \begin{column}{0.56\textwidth}
            \centering
            \begin{tikzpicture}[x=2.0cm, y=2.0cm, font=\scriptsize]
                \begin{scope}
                    \clip (-3.75,-2.55) rectangle (-1.15,-0.75);
                    \momValle
                    \draw[NaranjaBase!45, very thick] (-3.6,-1.2) -- (-2.406,-2.004);
                    % paso de GD desde x(2)
                    \draw[GrisTexto!70, thick, dashed, -stealth] (-2.406,-2.004) -- (-2.404,-0.99);
                    % inercia + gradiente
                    \draw[CelesteBrillante, very thick, -stealth] (-2.406,-2.004) -- (-1.83,-2.39);
                    \draw[NaranjaBase, very thick, -stealth] (-1.809,-2.406) -- (-1.807,-1.39);
                    \draw[VerdeNeon, ultra thick, -stealth] (-2.406,-2.004) -- (-1.83,-1.385);
                \end{scope}
                \draw[gray!50] (-3.75,-2.55) rectangle (-1.15,-0.75);
                \fill[GrisTexto!70] (-3.6,-1.2) circle (1.6pt) node[above right, inner sep=2pt] {$x^{(1)}$};
                \fill[GrisTexto] (-2.406,-2.004) circle (1.8pt) node[left, inner sep=3pt] {$x^{(2)}$};
                \fill[VerdeNeon] (-1.807,-1.360) circle (1.8pt) node[right, inner sep=3pt, text=GrisTexto] {$x^{(3)}$};
                \draw[GrisTexto, thick] (-2.404,-0.958) circle (2pt) node[left, inner sep=3pt] {GD};
                \node[CelesteBrillante, anchor=north east] at (-2.12,-2.25) {$\beta v^{(2)}$};
                \node[NaranjaBase, anchor=west] at (-1.78,-1.95) {$-\alpha g^{(2)}$};
                \node[VerdeNeon, anchor=south east, inner sep=1pt] at (-2.13,-1.62) {$v^{(3)}$};
            \end{tikzpicture}
        \end{column}
        \begin{column}{0.42\textwidth}
            \begin{block}{Iteración 2}
                \footnotesize
                $v^{(3)}=\textcolor{CelesteBrillante}{\beta v^{(2)}}\textcolor{NaranjaBase}{-\alpha g^{(2)}}$
            \end{block}
            \small
            \begin{itemize}
                \item \textcolor{CelesteBrillante}{Inercia}: sigue el paso anterior
                \item \textcolor{NaranjaBase}{Gradiente}: frena la oscilación
                \item \textcolor{VerdeNeon}{Resultado}: avanza \textbf{por} el valle
            \end{itemize}
            \vspace{0.3em}
            {\scriptsize GD (punteado) solo rebota de un lado al otro.\par}
        \end{column}
    \end{columns}
    \reflibro{Ec.~(5.20)--(5.21).}
\end{frame}

\begin{frame}{Momentum: Explicación Gráfica (3/3)}
    \begin{columns}[T]
        \begin{column}{0.56\textwidth}
            \centering
            \begin{tikzpicture}[x=1.15cm, y=1.15cm, font=\scriptsize]
                \begin{scope}
                    \clip (-3.9,-2.6) rectangle (0.6,0.5);
                    \momValle
                    \draw[CelesteBrillante, thick] plot coordinates {(-3.600,-1.200) (-2.406,-2.004) (-2.404,-0.958) (-1.710,-1.289) (-1.619,-0.724) (-1.204,-0.843) (-1.097,-0.530) (-0.842,-0.559) (-0.747,-0.381) (-0.587,-0.374) (-0.510,-0.270) (-0.407,-0.252)};
                    \draw[NaranjaBase, very thick] plot coordinates {(-3.600,-1.200) (-2.406,-2.004) (-1.807,-1.360) (-1.410,-0.444) (-0.726,-0.335) (-0.205,-0.321) (-0.044,-0.027) (0.043,0.126) (0.147,0.065) (0.161,0.046) (0.109,0.082) (0.078,0.064)};
                    \foreach \p in {(-2.406,-2.004), (-1.807,-1.360), (-1.410,-0.444), (-0.726,-0.335), (-0.205,-0.321), (-0.044,-0.027)}
                        \fill[NaranjaBase] \p circle (1.5pt);
                    \foreach \p in {(-2.404,-0.958), (-1.710,-1.289), (-1.619,-0.724), (-1.204,-0.843), (-1.097,-0.530)}
                        \fill[CelesteBrillante] \p circle (1.3pt);
                \end{scope}
                \draw[gray!50] (-3.9,-2.6) rectangle (0.6,0.5);
                \fill[white] (-3.6,-1.2) circle (1.6pt) node[above, inner sep=2pt, text=GrisTexto] {$x^{(1)}$};
                \fill[white] (0,0) circle (1.5pt) node[above left, inner sep=1pt] {$x^*$};
                \draw[NaranjaBase, very thick] (-3.8,-2.95) -- (-3.4,-2.95) node[right, text=GrisTexto] {Momentum};
                \draw[CelesteBrillante, thick] (-1.3,-2.95) -- (-0.9,-2.95) node[right, text=GrisTexto] {GD};
            \end{tikzpicture}
        \end{column}
        \begin{column}{0.42\textwidth}
            \begin{block}{Iteraciones 3, 4, \dots}
                \footnotesize
                La velocidad a lo largo del valle se \textbf{acumula}
            \end{block}
            \centering
            {\scriptsize Avance a lo largo del valle por iteración\par}
            \vspace{0.2em}
            \begin{tikzpicture}[x=0.55cm, y=2.0cm, font=\tiny]
                \draw[gray!60] (0.4,0) -- (7.6,0);
                \foreach \k/\m/\g in {1/0.632/0.632, 2/0.841/0.525, 3/0.802/0.435, 4/0.646/0.361, 5/0.459/0.300, 6/0.287/0.249, 7/0.152/0.207}{
                    \fill[NaranjaBase] ({\k-0.32},0) rectangle ({\k},\m);
                    \fill[CelesteBrillante] ({\k},0) rectangle ({\k+0.32},\g);
                    \node[below] at (\k,0) {$\k$};
                }
                \node[below] at (4,-0.13) {iteración $k$};
            \end{tikzpicture}

            \raggedright
            \vspace{0.3em}
            {\small $f(x^{(7)})=0.001$ \ (GD: $0.78$)\par}
        \end{column}
    \end{columns}
    \reflibro{Fig.~5.5 (mismo efecto en Rosenbrock); Ec.~(5.20)--(5.21).}
\end{frame}

% ---------------------------------------------------------------
% 2. Pseudocódigo
% ---------------------------------------------------------------
\begin{frame}{Momentum: Pseudocódigo}
    \scriptsize
    \begin{algorithmic}[1]
        \Require $\nabla f$, $x^{(1)}$, $\alpha>0$, $\beta\in[0,1)$, tolerancia $\epsilon$, $k_{\max}$
        \State $v \gets 0$
        \For{$k = 1, \dots, k_{\max}$}
            \State $g \gets \nabla f(x)$
            \State \textbf{if} $\lVert g \rVert < \epsilon$ \textbf{then break}
            \State $v \gets \beta\, v - \alpha\, g$ \Comment{(5.20)}
            \State $x \gets x + v$ \Comment{(5.21)}
        \EndFor
        \State \Return $x$
    \end{algorithmic}

    \vspace{1em}
    \small
    \begin{itemize}
        \item $\beta=0$ $\Rightarrow$ descenso de gradiente
        \item Típico: $\beta\approx0.9$
    \end{itemize}
    \reflibro{Algoritmo~5.3, Ec.~(5.20)--(5.21), Ejercicio~5.10.}
\end{frame}

% ---------------------------------------------------------------
% 3. Ejemplo paso a paso
% ---------------------------------------------------------------
\begin{frame}{Momentum: Ejemplo Paso a Paso (1/3)}
    \begin{exampleblock}{Planteamiento}
        \footnotesize
        $f(x)=\tfrac12x_1^2+5x_2^2$, \quad $\nabla f=(x_1,\ 10x_2)$, \quad
        $x^{(1)}=(-4,\,1)$, \quad $\alpha=0.17$, \ $\beta=0.5$
    \end{exampleblock}

    \begin{columns}[T]
        \begin{column}{0.5\textwidth}
            \begin{block}{Inicialización}
                \footnotesize
                $v=(0,\ 0)$
            \end{block}
            \begin{block}{Iteración $k=1$}
                \footnotesize
                $g^{(1)}=(-4,\ 10)$\\[3pt]
                $v^{(2)}=-0.17\,g^{(1)}=(0.68,\ -1.70)$\\[3pt]
                $x^{(2)}=(-3.32,\ -0.70)$
            \end{block}
            {\small Primer paso $=$ paso de GD\par}
        \end{column}
        \begin{column}{0.48\textwidth}
            \centering
            \begin{tikzpicture}[x=1.05cm, y=1.05cm, font=\scriptsize]
                \begin{scope}
                    \clip (-4.4,-1.3) rectangle (0.6,1.35);
                    \momElipses
                    \draw[NaranjaBase, very thick, -stealth] (-4,1) -- (-3.34,-0.65);
                \end{scope}
                \draw[gray!50] (-4.4,-1.3) rectangle (0.6,1.35);
                \fill[GrisTexto] (-4,1) circle (1.8pt) node[right, inner sep=3pt] {$x^{(1)}$};
                \fill[GrisTexto] (-3.32,-0.70) circle (1.8pt) node[right, inner sep=3pt] {$x^{(2)}$};
                \fill[white] (0,0) circle (1.6pt) node[above, inner sep=2pt] {$x^*$};
                \node[NaranjaBase, anchor=east] at (-3.72,0.05) {$v^{(2)}$};
            \end{tikzpicture}
        \end{column}
    \end{columns}
\end{frame}

\begin{frame}{Momentum: Ejemplo Paso a Paso (2/3)}
    \begin{columns}[T]
        \begin{column}{0.56\textwidth}
            \begin{block}{Iteración $k=2$}
                \footnotesize
                $g^{(2)}=(-3.32,\ -7.00)$\\[3pt]
                $\textcolor{CelesteBrillante}{\beta v^{(2)}}=(0.34,\ -0.85)$\\[3pt]
                $\textcolor{NaranjaBase}{-\alpha g^{(2)}}=(0.56,\ 1.19)$\\[3pt]
                $\textcolor{VerdeNeon}{v^{(3)}}=(0.90,\ 0.34)$\\[3pt]
                $x^{(3)}=(-2.42,\ -0.36)$
            \end{block}
            \small
            \begin{itemize}
                \item $x_1$: velocidad \textbf{crece}
                \item $x_2$: oscilación \textbf{se cancela}
            \end{itemize}
        \end{column}
        \begin{column}{0.42\textwidth}
            \centering
            \begin{tikzpicture}[x=1.45cm, y=1.45cm, font=\scriptsize]
                \begin{scope}
                    \clip (-4.25,-1.75) rectangle (-1.9,1.15);
                    \momElipses
                    \draw[GrisTexto!60, thick, -stealth] (-4,1) -- (-3.34,-0.65);
                    % paso de GD desde x(2)
                    \draw[GrisTexto!70, thick, dashed, -stealth] (-3.32,-0.70) -- (-2.77,0.46);
                    % momentum: inercia + gradiente
                    \draw[CelesteBrillante, very thick, -stealth] (-3.32,-0.70) -- (-2.99,-1.53);
                    \draw[NaranjaBase, very thick, -stealth] (-2.98,-1.55) -- (-2.43,-0.39);
                    \draw[VerdeNeon, ultra thick, -stealth] (-3.32,-0.70) -- (-2.45,-0.37);
                \end{scope}
                \draw[gray!50] (-4.25,-1.75) rectangle (-1.9,1.15);
                \fill[GrisTexto] (-4,1) circle (1.6pt) node[right, inner sep=3pt] {$x^{(1)}$};
                \fill[GrisTexto] (-3.32,-0.70) circle (1.6pt) node[left, inner sep=2pt] {$x^{(2)}$};
                \fill[VerdeNeon] (-2.4156,-0.36) circle (1.8pt) node[right, inner sep=2pt, text=GrisTexto] {$x^{(3)}$};
                \draw[GrisTexto, thick] (-2.7556,0.49) circle (2pt);
                \node[anchor=west, text=GrisTexto] at (-2.68,0.52) {GD};
                \node[anchor=east, text=CelesteBrillante] at (-3.2,-1.25) {$\beta v^{(2)}$};
                \node[anchor=west, text=NaranjaBase] at (-2.66,-1.05) {$-\alpha g^{(2)}$};
                \node[anchor=south, text=VerdeNeon] at (-2.75,-0.42) {$v^{(3)}$};
            \end{tikzpicture}
        \end{column}
    \end{columns}
\end{frame}

\begin{frame}{Momentum: Ejemplo Paso a Paso (3/3)}
    \begin{columns}[T]
        \begin{column}{0.56\textwidth}
            \begin{block}{Iteración $k=3$}
                \footnotesize
                $g^{(3)}=(-2.42,\ -3.60)$\\[3pt]
                $v^{(4)}=(0.86,\ 0.78)$\\[3pt]
                $x^{(4)}=(-1.55,\ 0.42)$
            \end{block}
            \begin{exampleblock}{Resultado ($\epsilon=10^{-6}$)}
                \footnotesize
                Momentum: $k=49$ \quad GD: $k=83$
            \end{exampleblock}
        \end{column}
        \begin{column}{0.42\textwidth}

            \centering
            \begin{tikzpicture}[x=0.95cm, y=0.95cm, font=\scriptsize]
                \begin{scope}
                    \clip (-4.4,-1.15) rectangle (0.5,1.3);
                    \momElipses
                    \draw[CelesteBrillante, thick] plot coordinates {(-4.000,1.000) (-3.320,-0.700) (-2.756,0.490) (-2.287,-0.343) (-1.898,0.240) (-1.576,-0.168) (-1.308,0.118) (-1.085,-0.082) (-0.901,0.058) (-0.748,-0.040) (-0.621,0.028) (-0.515,-0.020) (-0.428,0.014) (-0.355,-0.010) (-0.295,0.007) (-0.244,-0.005) (-0.203,0.003) (-0.168,-0.002) (-0.140,0.002) (-0.116,-0.001) (-0.096,0.001)};
                    \draw[NaranjaBase, very thick] plot coordinates {(-4.000,1.000) (-3.320,-0.700) (-2.416,-0.360) (-1.553,0.422) (-0.857,0.096) (-0.364,-0.230) (-0.055,-0.002) (0.108,0.115) (0.172,-0.022) (0.174,-0.053) (0.146,0.022) (0.107,0.022) (0.069,-0.015) (0.039,-0.008) (0.017,0.009) (0.003,0.002) (-0.004,-0.005) (-0.007,-0.000) (-0.008,0.003) (-0.006,-0.000) (-0.005,-0.001)};
                    \foreach \p in {(-3.320,-0.700), (-2.416,-0.360), (-1.553,0.422), (-0.857,0.096), (-0.364,-0.230), (-0.055,-0.002)}
                        \fill[NaranjaBase] \p circle (1.4pt);
                    \foreach \p in {(-2.756,0.490), (-2.287,-0.343), (-1.898,0.240), (-1.576,-0.168), (-1.308,0.118)}
                        \fill[CelesteBrillante] \p circle (1.2pt);
                \end{scope}
                \draw[gray!50] (-4.4,-1.15) rectangle (0.5,1.3);
                \fill[white] (-4,1) circle (1.6pt) node[right, inner sep=3pt, text=GrisTexto] {$x^{(1)}$};
                \fill[white] (0,0) circle (1.4pt);
                \draw[NaranjaBase, very thick] (-4.3,-1.45) -- (-3.9,-1.45) node[right, text=GrisTexto] {Momentum};
                \draw[CelesteBrillante, thick] (-1.6,-1.45) -- (-1.2,-1.45) node[right, text=GrisTexto] {GD};
            \end{tikzpicture}

            {\tiny Primeras 20 iteraciones de cada método.\par}
        \end{column}
    \end{columns}

    \vspace{0.3em}
    \centering
    {\scriptsize
    \begin{tabular}{l c c c c c c c}
        \toprule
        $k$ & 1 & 2 & 3 & 4 & 5 & 6 & 7 \\
        \midrule
        $f(x^{(k)})$ Momentum & 13 & 7.961 & \textbf{3.566} & \textbf{2.096} & \textbf{0.413} & \textbf{0.331} & \textbf{0.002} \\
        $f(x^{(k)})$ GD ($\beta=0$) & 13 & 7.961 & 4.997 & 3.204 & 2.090 & 1.383 & 0.924 \\
        \bottomrule
    \end{tabular}\par}
\end{frame}

% ---------------------------------------------------------------
% 4. Código
% ---------------------------------------------------------------
\begin{frame}[fragile]{Momentum: Código}
\begin{lstlisting}[language=Python]
import numpy as np

def momentum(grad, x, alpha=0.01, beta=0.9, k_max=1000, eps=1e-6):
    v = np.zeros_like(x)                  # velocidad inicial v = 0
    for k in range(1, k_max + 1):
        g = grad(x)
        if np.linalg.norm(g) < eps:       # criterio de parada
            break
        v = beta*v - alpha*g              # (5.20) inercia + gradiente
        x = x + v                         # (5.21) avanzar con la velocidad
    return x, k

# Ejemplo: f(x) = 0.5*x1^2 + 5*x2^2   ->   grad f = (x1, 10*x2)
grad = lambda x: np.array([x[0], 10*x[1]])
x0 = np.array([-4.0, 1.0])
momentum(grad, x0, alpha=0.17, beta=0.5, k_max=3)  # x(4) = [-1.5527, 0.422]
momentum(grad, x0, alpha=0.17, beta=0.5)           # -> [0, 0] en k = 49
momentum(grad, x0, alpha=0.17, beta=0.0)           # GD: k = 83
\end{lstlisting}
\end{frame}

% ---------------------------------------------------------------
% 5. Análisis de convergencia
% ---------------------------------------------------------------
\begin{frame}{Momentum: Análisis de Convergencia (1/2)}
    \begin{columns}[T]
        \begin{column}{0.5\textwidth}
            \centering
            \begin{tikzpicture}[x=0.085cm, y=0.26cm, font=\tiny]
                \draw[gray!60] (0,-12) -- (50,-12);
                \draw[gray!60] (0,-12) -- (0,2);
                \foreach \yy/\lb in {0/{10^{0}}, -4/{10^{-4}}, -8/{10^{-8}}, -12/{10^{-12}}}{
                    \draw[gray!60] (-0.8,\yy) -- (0,\yy);
                    \draw[gray!25] (0,\yy) -- (50,\yy);
                    \node[left] at (-0.8,\yy) {$\lb$};
                }
                \foreach \xx in {0,10,...,50}
                    \draw[gray!60] (\xx,-12) -- (\xx,-12.4) node[below] {$\xx$};
                \node[below] at (25,-13.3) {iteración $k$};
                \node[rotate=90] at (-11,-5) {$f(x^{(k)})-f^*$};
                \draw[CelesteBrillante, thick] plot coordinates {(0,1.11) (1,0.90) (2,0.70) (3,0.51) (4,0.32) (5,0.14) (6,-0.03) (7,-0.21) (8,-0.37) (9,-0.54) (10,-0.71) (11,-0.87) (12,-1.03) (13,-1.20) (14,-1.36) (15,-1.52) (16,-1.69) (17,-1.85) (18,-2.01) (19,-2.17) (20,-2.33) (21,-2.50) (22,-2.66) (23,-2.82) (24,-2.98) (25,-3.14) (26,-3.30) (27,-3.47) (28,-3.63) (29,-3.79) (30,-3.95) (31,-4.11) (32,-4.28) (33,-4.44) (34,-4.60) (35,-4.76) (36,-4.92) (37,-5.09) (38,-5.25) (39,-5.41) (40,-5.57) (41,-5.73) (42,-5.89) (43,-6.06) (44,-6.22) (45,-6.38) (46,-6.54) (47,-6.70) (48,-6.87) (49,-7.03) (50,-7.19)};
                \draw[RojoAlerta, thick] plot coordinates {(0,1.11) (1,0.90) (2,0.89) (3,0.06) (4,0.74) (5,0.26) (6,0.89) (7,0.58) (8,0.83) (9,0.27) (10,0.48) (11,-0.28) (12,0.34) (13,0.31) (14,0.43) (15,0.43) (16,0.19) (17,0.18) (18,-0.70) (19,0.06) (20,-0.36) (21,0.22) (22,-0.10) (23,0.13) (24,-0.46) (25,-0.22) (26,-0.92) (27,-0.33) (28,-0.34) (29,-0.25) (30,-0.26) (31,-0.53) (32,-0.52) (33,-1.46) (34,-0.60) (35,-0.99) (36,-0.46) (37,-0.79) (38,-0.57) (39,-1.20) (40,-0.93) (41,-1.56) (42,-1.00) (43,-1.00) (44,-0.93) (45,-0.95) (46,-1.25) (47,-1.22) (48,-2.20) (49,-1.27) (50,-1.63)};
                \draw[NaranjaBase, very thick] plot coordinates {(0,1.11) (1,0.90) (2,0.55) (3,0.32) (4,-0.38) (5,-0.48) (6,-2.81) (7,-1.14) (8,-1.76) (9,-1.53) (10,-1.89) (11,-2.09) (12,-2.45) (13,-2.97) (14,-3.24) (15,-4.55) (16,-3.86) (17,-4.57) (18,-4.21) (19,-4.66) (20,-4.74) (21,-5.23) (22,-5.56) (23,-6.07) (24,-6.75) (25,-6.65) (26,-7.21) (27,-6.93) (28,-7.40) (29,-7.42) (30,-8.00) (31,-8.19) (32,-8.94) (33,-9.19) (34,-9.53) (35,-9.76) (36,-9.69) (37,-10.08) (38,-10.13) (39,-10.73) (40,-10.85) (41,-11.81) (42,-11.75) (43,-12.00) (44,-12.00) (45,-12.00) (46,-12.00) (47,-12.00) (48,-12.00) (49,-12.00) (50,-12.00)};
                \node[anchor=west, text=CelesteBrillante] at (30,-3.1) {$\beta=0$ (GD)};
                \node[anchor=west, text=RojoAlerta] at (36,1.0) {$\beta=0.9$};
                \node[anchor=west, text=NaranjaBase] at (21,-9.6) {$\beta=0.5$};
            \end{tikzpicture}

            {\tiny Ejemplo anterior, variando solo $\beta$.\par}
        \end{column}
        \begin{column}{0.48\textwidth}
            \begin{alertblock}{Convergencia lineal: error $\sim\rho^k$}
                \footnotesize
                $\beta$ moderado acelera; excesivo oscila.
            \end{alertblock}
            \centering
            \small
            \begin{tabular}{l c c}
                \toprule
                $\beta$ & $\rho$ & iteraciones \\
                \midrule
                $0$ (GD) & $0.83$ & $83$ \\
                $0.5$ & $0.71$ & $\mathbf{49}$ \\
                $0.9$ & $0.95$ & $289$ \\
                \bottomrule
            \end{tabular}
        \end{column}
    \end{columns}
    \reflibro{Ec.~(5.20)--(5.21); experimento propio con el ejemplo del paso a paso.}
\end{frame}

\begin{frame}{Momentum: Análisis de Convergencia (2/2)}
    \begin{columns}[T]
        \begin{column}{0.5\textwidth}
            \centering
            {\scriptsize\textbf{Región de convergencia en una cuadrática}\par}
            \vspace{0.2em}
            \begin{tikzpicture}[x=1.15cm, y=3.0cm, font=\tiny]
                % Región estable: 0 <= beta < 1, 0 < alpha*lambda < 2(1+beta)
                \fill[VerdeNeon!35!black] (0,0) -- (2,0) -- (4,1) -- (0,1) -- cycle;
                % Subregión con raíces complejas (convergencia oscilatoria)
                \fill[VerdeNeon!16!black]
                    plot[domain=0:1, samples=40, variable=\t] ({(1-sqrt(\t))^2}, \t)
                    -- plot[domain=1:0, samples=40, variable=\t] ({(1+sqrt(\t))^2}, \t) -- cycle;
                \draw[NaranjaBase!80, domain=0:1, samples=40, variable=\t] plot ({(1-sqrt(\t))^2}, \t);
                \draw[NaranjaBase!80, domain=0:1, samples=40, variable=\t] plot ({(1+sqrt(\t))^2}, \t);
                \draw[VerdeNeon, thick] (2,0) -- (4,1);
                \draw[gray!60] (0,0) -- (4.2,0);
                \draw[gray!60] (0,0) -- (0,1.08);
                \draw[gray!40, dashed] (0,1) -- (4,1);
                \foreach \xx in {0,1,2,3,4}
                    \draw[gray!60] (\xx,0) -- (\xx,-0.03) node[below] {$\xx$};
                \foreach \yy in {0,0.5,1}
                    \draw[gray!60] (0,\yy) -- (-0.08,\yy) node[left] {$\yy$};
                \node[below] at (2.1,-0.12) {$\alpha\lambda$ \ ($\lambda$: autovalor de la Hessiana)};
                \node at (-0.42,0.75) {$\beta$};
                \node[RojoAlerta] at (3.45,0.34) {diverge};
                \node[RojoAlerta] at (3.45,0.25) {$\alpha\lambda>2(1+\beta)$};
                \node[NaranjaBase] at (1.55,0.8) {converge oscilando};
                \node[NaranjaBase] at (1.55,0.72) {$|r|=\sqrt{\beta}$};
                % Leyenda
                \fill[VerdeNeon!35!black] (0,-0.36) rectangle (0.2,-0.30);
                \node[anchor=west] at (0.22,-0.33) {converge monótono};
                \fill[VerdeNeon!16!black] (2.1,-0.36) rectangle (2.3,-0.30);
                \draw[NaranjaBase!80] (2.1,-0.36) rectangle (2.3,-0.30);
                \node[anchor=west] at (2.32,-0.33) {converge oscilando};
                % puntos del ejemplo
                \fill[white] (1.7,0.5) circle (1.5pt);
                \fill[white] (0.17,0.5) circle (1.5pt);
                \node[anchor=north, text=white] at (0.95,0.5) {ejemplo};
                \fill[RojoAlerta] (2,0) circle (1.8pt);
                \fill[VerdeNeon] (2,0.5) circle (1.6pt);
                \node[anchor=south west, text=RojoAlerta, inner sep=1pt] at (2.05,0.02) {GD, $\alpha=0.2$};
                \node[anchor=north, text=VerdeNeon, inner sep=2pt] at (2.05,0.49) {Mom., $\alpha=0.2$};
            \end{tikzpicture}
        \end{column}
        \begin{column}{0.48\textwidth}
            \small
            \begin{itemize}
                \item Converge si $\alpha\lambda<2(1+\beta)$\\
                      (GD: $\alpha\lambda<2$) $\Rightarrow$ \textbf{pasos más grandes}
                \item Razón óptima:\\[2pt]
                      $\frac{\sqrt{\kappa}-1}{\sqrt{\kappa}+1}$ vs.\ $\frac{\kappa-1}{\kappa+1}$ (GD)\\[2pt]
                      $\kappa=10$: $0.52$ vs.\ $0.82$
                \item[$-$] Sensible a $\alpha,\beta$; sobrepasa el fondo (Nesterov)
            \end{itemize}
        \end{column}
    \end{columns}
    \reflibro{Ejercicio~5.10, §5.4. Región y razón óptima: resultados clásicos (Polyak, 1964), no del libro.}
\end{frame}

\subsection{Adagrad}

\begin{frame}{Adagrad: Explicación Gráfica}
    \begin{block}{Concepto principal}
        [Explica la adaptación de la tasa de aprendizaje por dimensión según frecuencias de actualización].
    \end{block}
    \begin{itemize}
        \item {[Idea gráfica 1: Escalamiento elipsoidal vs pasos simétricos estándar].}
        \item {[Idea gráfica 2: Atenuación progresiva del paso a lo largo de las iteraciones].}
    \end{itemize}
\end{frame}

\begin{frame}{Adagrad: Pseudocódigo}
    \begin{algorithmic}[1]
        \State \textbf{Input:} {[Tasa base \(\alpha\), parámetro de suavizado \(\epsilon\), punto inicial \(x\)]}
        \State Inicializar acumulador de gradientes al cuadrado \(s \leftarrow 0\)
        \While{[Condición de parada]}
            \State Calcular gradiente actual \(g\)
            \State Acumular cuadrados de componentes: \(s \leftarrow s + g \odot g\)
            \State Actualizar variables escalando por \(1 / \sqrt{s + \epsilon}\)
        \EndWhile
        \State \textbf{Return} {[Punto óptimo]}
    \end{algorithmic}
\end{frame}

\begin{frame}{Adagrad: Ejemplo Paso a Paso}
    \begin{exampleblock}{Planteamiento del problema}
        [Define función con diferentes escalas por eje y un punto de partida].
    \end{exampleblock}
    \begin{itemize}
        \item \textbf{Paso 1:} {[Cálculo del gradiente y primer vector acumulador \(s\)].}
        \item \textbf{Paso 2:} {[Actualización diferenciada para cada componente].}
        \item \textbf{Paso 3:} {[Resultado tras las primeras iteraciones].}
    \end{itemize}
\end{frame}

\begin{frame}[fragile]{Adagrad: Código}
\begin{lstlisting}[language=Python]
def adagrad(grad_f, x0, alpha=0.1, eps=1e-8, max_iter=100):
    # TODO: Inicializar s acumulador en ceros
    # for t in range(max_iter):
    #     g = grad_f(x)
    #     s += g**2
    #     x -= (alpha / (np.sqrt(s) + eps)) * g
    pass
\end{lstlisting}
\end{frame}

\begin{frame}{Adagrad: Análisis de Convergencia}
    \begin{alertblock}{Propiedades de Convergencia}
        [Discute cómo decae la tasa de aprendizaje efectiva y su impacto final].
    \end{alertblock}
    \begin{itemize}
        \item {[Idoneidad para problemas con características o gradientes dispersos (sparse)].}
        \item {[Limitación principal: congelamiento prematuro del paso].}
    \end{itemize}
\end{frame}
\subsection{RMSProp}

% =========================================================
% 1. EXPLICACIÓN GRÁFICA
% =========================================================
\begin{frame}{RMSProp: Explicación Gráfica}

\footnotesize

\begin{columns}[T,onlytextwidth]

\column{0.52\textwidth}

\begin{center}
\begin{tikzpicture}[scale=0.58]

    \draw[->, GrisTexto] (-2.6,0) -- (2.8,0) node[right] {$x$};
    \draw[->, GrisTexto] (0,-2.0) -- (0,2.2) node[above] {$y$};

    \draw[CelesteBrillante, thick] (0,0) ellipse (2.1 and 1.25);
    \draw[CelesteBrillante, thick] (0,0) ellipse (1.55 and 0.92);
    \draw[CelesteBrillante, thick] (0,0) ellipse (1.0 and 0.58);
    \draw[CelesteBrillante, thick] (0,0) ellipse (0.45 and 0.26);

    \fill[NaranjaBase] (2.0,1.35) circle (1.7pt);
    \fill[NaranjaBase] (1.35,0.72) circle (1.7pt);
    \fill[NaranjaBase] (0.88,0.36) circle (1.7pt);
    \fill[NaranjaBase] (0.42,0.14) circle (1.7pt);
    \fill[VerdeNeon] (0.10,0.03) circle (2pt);

    \draw[NaranjaBase, thick, ->] (2.0,1.35) -- (1.35,0.72);
    \draw[NaranjaBase, thick, ->] (1.35,0.72) -- (0.88,0.36);
    \draw[NaranjaBase, thick, ->] (0.88,0.36) -- (0.42,0.14);
    \draw[NaranjaBase, thick, ->] (0.42,0.14) -- (0.10,0.03);

\end{tikzpicture}

\textcolor{VerdeNeon}{Trayectoria hacia el mínimo}
\end{center}

\column{0.45\textwidth}

\textbf{¿Qué está ocurriendo?}

\begin{itemize}
    \item RMSProp guarda un promedio reciente de $g^2$.
    \item Gradiente grande $\Rightarrow$ paso menor.
    \item Gradiente pequeño $\Rightarrow$ paso relativamente mayor.
    \item Los gradientes antiguos pierden peso.
\end{itemize}

\[
s_k=\gamma s_{k-1}+(1-\gamma)g_k^2
\]

\[
x_{k+1}
=
x_k-
\frac{\alpha g_k}{\sqrt{s_k}+\epsilon}
\]

\end{columns}

\end{frame}

% =========================================================
% 2. PSEUDOCÓDIGO
% =========================================================
\begin{frame}{RMSProp: Pseudocódigo}

\small

\begin{algorithmic}[1]

\Require $x,\alpha,\gamma,\epsilon$

\State $s\gets0$

\While{no se cumpla la condición de parada}

    \State $g\gets\nabla f(x)$

    \State
    $s\gets
    \gamma s+(1-\gamma)(g\odot g)$

    \State
    $x\gets
    x-\dfrac{\alpha g}{\sqrt{s}+\epsilon}$

\EndWhile

\State \Return $x$

\end{algorithmic}

\vspace{0.35cm}

\begin{block}{Interpretación}
Un gradiente grande aumenta $s$ y reduce el paso.
Un gradiente pequeño permite un movimiento relativamente mayor.
\end{block}

\begin{center}
\footnotesize
$\gamma=0.9$:
90\% historial + 10\% gradiente actual.
\end{center}

\end{frame}


% =========================================================
% 3. EJEMPLO
% =========================================================
\begin{frame}{RMSProp: Ejemplo Paso a Paso}

\small

\[
f(x)=\frac12x^2,
\qquad
f'(x)=x
\]

\[
x_0=2,\quad
\alpha=0.1,\quad
\gamma=0.9,\quad
\epsilon=0.01,\quad
s_0=0
\]

\vspace{0.15cm}

\begin{center}
\scriptsize
\begin{tabular}{c|c|c|c|c}
\toprule
$k$ & $x_k$ & $g_k$ & $s_{k+1}$ & $x_{k+1}$ \\
\midrule
0 & 2.0000 & 2.0000 & 0.4000 & 1.6887 \\
1 & 1.6887 & 1.6887 & 0.6452 & 1.4810 \\
2 & 1.4810 & 1.4810 & 0.8000 & 1.3173 \\
\bottomrule
\end{tabular}
\end{center}

\vspace{0.2cm}

\footnotesize
\textbf{Primera iteración:}

\[
g_0=2
\qquad
s_1=0.9(0)+0.1(2^2)=0.4
\]

\[
x_1
=
2-\frac{0.1(2)}
{\sqrt{0.4}+0.01}
\approx1.6887
\]

\begin{center}
\textcolor{VerdeNeon}{
$2\rightarrow1.6887\rightarrow1.4810\rightarrow1.3173$
}
\end{center}

\end{frame}


% =========================================================
% 4. CÓDIGO
% =========================================================
\begin{frame}[fragile]{RMSProp: Código}

\begin{lstlisting}[
language=Python,
basicstyle=\ttfamily\scriptsize,
showstringspaces=false
]
def rmsprop(grad, x, alpha=.1,
            gamma=.9, eps=.01, n=10):
    s = 0.0

    for _ in range(n):
        g = grad(x)
        s = gamma*s + (1-gamma)*g**2
        x -= alpha*g/(s**0.5 + eps)

    return x


def grad(x):
    return x


print(rmsprop(grad, 2.0))
\end{lstlisting}

\end{frame}
\subsection{Adadelta}

% =========================================================
% 1. EXPLICACIÓN GRÁFICA
% =========================================================
\begin{frame}{Adadelta: Explicación Gráfica}

\footnotesize

\begin{center}

\begin{tikzpicture}[scale=0.82]

    \node[
        draw=CelesteBrillante,
        rounded corners,
        text=GrisTexto,
        minimum width=2.6cm,
        minimum height=0.7cm
    ] at (0,1.7) (g) {Gradiente $g_k$};

    \node[
        draw=NaranjaBase,
        rounded corners,
        text=GrisTexto,
        align=center,
        minimum width=3.2cm,
        minimum height=0.85cm
    ] at (-3.1,0) (s) {Promedio de\\$g^2$: $s_k$};

    \node[
        draw=VerdeNeon,
        rounded corners,
        text=GrisTexto,
        align=center,
        minimum width=3.2cm,
        minimum height=0.85cm
    ] at (3.1,0) (u) {Promedio de\\$(\Delta x)^2$: $u_k$};

    \node[
        draw=GrisTexto,
        rounded corners,
        text=GrisTexto,
        minimum width=3.0cm,
        minimum height=0.75cm
    ] at (0,-1.7) (dx) {Actualización $\Delta x_k$};

    \draw[->, thick, GrisTexto] (g) -- (s);
    \draw[->, thick, GrisTexto] (g) -- (dx);
    \draw[->, thick, GrisTexto] (s) -- (dx);
    \draw[->, thick, GrisTexto] (u) -- (dx);

    \draw[->, thick, GrisTexto]
        (dx.east) -- (4.6,-1.7) -- (4.6,0) -- (u.east);

\end{tikzpicture}

\end{center}

\vspace{0.1cm}

\[
\Delta x_k
=
-\frac{\sqrt{u_{k-1}}+\epsilon}
{\sqrt{s_k}+\epsilon}\,g_k
\]

\begin{center}
\textbf{Gradientes} $\rightarrow s$
\qquad
\textbf{actualizaciones} $\rightarrow u$

\vspace{0.1cm}

\textcolor{VerdeNeon}{
\textbf{Adadelta usa ambas memorias y no necesita $\alpha$.}
}
\end{center}

\end{frame}


% =========================================================
% 2. PSEUDOCÓDIGO
% =========================================================
\begin{frame}{Adadelta: Pseudocódigo}

\footnotesize

\begin{algorithmic}[1]

\Require $x,\gamma_s,\gamma_x,\epsilon$

\State $s\gets0$
\State $u\gets0$

\While{no se cumpla la condición de parada}

    \State $g\gets\nabla f(x)$

    \State
    $s\gets
    \gamma_s s+(1-\gamma_s)(g\odot g)$

    \State
    $\Delta x\gets
    -\dfrac{\sqrt{u}+\epsilon}
    {\sqrt{s}+\epsilon}\odot g$

    \State
    $u\gets
    \gamma_xu+
    (1-\gamma_x)(\Delta x\odot\Delta x)$

    \State $x\gets x+\Delta x$

\EndWhile

\State \Return $x$

\end{algorithmic}

\vspace{0.25cm}

\begin{block}{Diferencia clave}
RMSProp todavía usa $\alpha$.
Adadelta usa $u$ para ajustar automáticamente
el tamaño de la actualización.
\end{block}

\end{frame}


% =========================================================
% 3. EJEMPLO
% =========================================================
\begin{frame}{Adadelta: Ejemplo Paso a Paso}

\small

\[
f(x)=\frac12x^2,
\qquad
f'(x)=x
\]

\[
x_0=2,\quad
\gamma_s=\gamma_x=0.9,\quad
\epsilon=0.1,\quad
s_0=u_0=0
\]

\vspace{0.15cm}

\begin{center}
\scriptsize
\begin{tabular}{c|c|c|c|c|c}
\toprule
$k$ & $x_k$ & $g_k$ & $s_{k+1}$ &
$\Delta x_k$ & $x_{k+1}$ \\
\midrule
0 & 2.0000 & 2.0000 & 0.4000 & -0.2731 & 1.7269 \\
1 & 1.7269 & 1.7269 & 0.6582 & -0.3531 & 1.3738 \\
2 & 1.3738 & 1.3738 & 0.7811 & -0.3330 & 1.0408 \\
\bottomrule
\end{tabular}
\end{center}

\vspace{0.15cm}

\footnotesize
\textbf{Primera iteración:}

\[
s_1=0.9(0)+0.1(2^2)=0.4
\]

\[
\Delta x_0=
-\frac{0.1}{\sqrt{0.4}+0.1}(2)
\approx-0.2731
\]

\[
x_1=2+\Delta x_0=1.7269
\]

\end{frame}


% =========================================================
% 4. CÓDIGO
% =========================================================
\begin{frame}[fragile]{Adadelta: Código}

\begin{lstlisting}[
language=Python,
basicstyle=\ttfamily\scriptsize,
showstringspaces=false
]
def adadelta(grad, x, gs=.9,
             gx=.9, eps=.1, n=10):
    s = 0.0
    u = 0.0

    for _ in range(n):
        g = grad(x)
        s = gs*s + (1-gs)*g**2

        dx = -((u**0.5 + eps) /
               (s**0.5 + eps))*g

        u = gx*u + (1-gx)*dx**2
        x += dx

    return x


def grad(x):
    return x


print(adadelta(grad, 2.0))
\end{lstlisting}

\end{frame}
\subsection{Adam}

\begin{frame}{Adam: Explicación Gráfica}
    \begin{block}{Concepto principal}
        [Explica la combinación de estimación de primer momento (inercia) y segundo momento (escala adaptativa)].
    \end{block}
    \begin{itemize}
        \item {[Idea gráfica 1: Corrección de sesgo durante las primeras iteraciones cerca del origen].}
        \item {[Idea gráfica 2: Trayectoria en superficies con diferente curvatura].}
    \end{itemize}
\end{frame}

\begin{frame}{Adam: Pseudocódigo}
    \begin{algorithmic}[1]
        \State \textbf{Input:} {[Parámetros \(\alpha, \beta_1, \beta_2, \epsilon\), punto inicial \(x\)]}
        \State Inicializar momentos \(m \leftarrow 0\), \(v \leftarrow 0\), paso \(t \leftarrow 0\)
        \While{[Condición de parada]}
            \State \(t \leftarrow t + 1\)
            \State Actualizar momentos sesgados \(m\) y \(v\) con el gradiente \(g\)
            \State Aplicar corrección de sesgo: \(\hat{m} \leftarrow m / (1-\beta_1^t)\), \(\hat{v} \leftarrow v / (1-\beta_2^t)\)
            \State Actualizar parámetros: \(x \leftarrow x - \alpha \hat{m} / (\sqrt{\hat{v}} + \epsilon)\)
        \EndWhile
        \State \textbf{Return} {[Punto óptimo]}
    \end{algorithmic}
\end{frame}

\begin{frame}{Adam: Ejemplo Paso a Paso}
    \begin{exampleblock}{Planteamiento del problema}
        [Define función cuadrática o de prueba junto con hiperparámetros estándar].
    \end{exampleblock}
    \begin{itemize}
        \item \textbf{Paso 1:} {[Cálculo de \(m_1, v_1\) y aplicación de corrección de sesgo].}
        \item \textbf{Paso 2:} {[Cálculo de la primera actualización de \(x\)].}
        \item \textbf{Paso 3:} {[Cálculos para la iteración \(t=2\)].}
    \end{itemize}
\end{frame}

\begin{frame}[fragile]{Adam: Código}
\begin{lstlisting}[language=Python]
def adam(grad_f, x0, alpha=0.001, beta1=0.9, beta2=0.999, eps=1e-8, max_iter=100):
    # TODO: Inicializar m, v y contador t
    # for t in range(1, max_iter + 1):
    #     g = grad_f(x)
    #     m = beta1 * m + (1 - beta1) * g
    #     v = beta2 * v + (1 - beta2) * (g ** 2)
    #     m_hat = m / (1 - beta1**t)
    #     v_hat = v / (1 - beta2**t)
    #     x = x - alpha * m_hat / (np.sqrt(v_hat) + eps)
    pass
\end{lstlisting}
\end{frame}

\begin{frame}{Adam: Análisis de Convergencia}
    \begin{alertblock}{Propiedades de Convergencia}
        [Resumir comportamiento general en problemas estocásticos y convexos].
    \end{alertblock}
    \begin{itemize}
        \item {[Impacto de la corrección de sesgo en los primeros pasos].}
        \item {[Discusión sobre casos donde puede no converger (variantes AMSGrad)].}
    \end{itemize}
\end{frame}
\subsection{Hypergradient Descent}

% ---------------------------------------------------------------
% 1. Explicación gráfica
% ---------------------------------------------------------------
\begin{frame}{Hypergradient Descent: Explicación Gráfica}
    \begin{block}{Concepto principal}
        \scriptsize
        La tasa de aprendizaje $\alpha$ es el hiperparámetro crítico de todo método de
        primer orden. La idea: tratarla como \emph{una variable más} y aplicarle
        \textbf{descenso de gradiente a ella misma}.
        \vspace{-0.3em}
        \[
            \frac{\partial f(\mathbf{x}^{(k)})}{\partial \alpha}
            = \nabla f(\mathbf{x}^{(k)})^{\top}
              \frac{\partial \mathbf{x}^{(k)}}{\partial \alpha}
            = -\,\nabla f(\mathbf{x}^{(k)})^{\top} \nabla f(\mathbf{x}^{(k-1)})
        \]
        \vspace{-0.5em}

        El \textbf{hipergradiente} es el producto punto de los \textbf{dos últimos
        gradientes}: información que ya se tiene, gratis.
    \end{block}

    \vspace{-0.2em}
    \begin{center}
    \begin{tikzpicture}[x=0.42cm, y=0.33cm, font=\tiny]
        % ---- panel izquierdo: gradientes alineados ----
        \draw[CelesteBrillante, thick, domain=-2.4:2.4, samples=40, smooth]
              plot (\x, {0.45*\x*\x});
        \fill[RojoAlerta] (2.2,2.178) circle (1.4pt);
        \fill[RojoAlerta] (1.4,0.882) circle (1.4pt);
        \draw[NaranjaBase, ->, very thick] (2.2,3.4) -- (1.0,3.4);
        \draw[NaranjaBase, ->, very thick] (1.4,4.4) -- (0.2,4.4);
        \node[above, text=VerdeNeon] at (0.9,4.8) {$g^{(k)}\!\cdot g^{(k-1)} > 0$};
        \node[below, text=GrisTexto, align=center] at (0,-0.4)
             {mismo sentido: voy bien\\ \textcolor{VerdeNeon}{$\alpha$ \textbf{aumenta}}};
        % ---- panel derecho: gradientes opuestos ----
        \begin{scope}[xshift=5.4cm]
        \draw[CelesteBrillante, thick, domain=-2.4:2.4, samples=40, smooth]
              plot (\x, {0.45*\x*\x});
        \fill[RojoAlerta] (2.2,2.178) circle (1.4pt);
        \fill[RojoAlerta] (-1.8,1.458) circle (1.4pt);
        \draw[NaranjaBase, ->, very thick] (2.2,3.4) -- (1.0,3.4);
        \draw[NaranjaBase, ->, very thick] (-1.8,4.4) -- (-0.6,4.4);
        \node[above, text=RojoAlerta] at (0.2,4.8) {$g^{(k)}\!\cdot g^{(k-1)} < 0$};
        \node[below, text=GrisTexto, align=center] at (0,-0.4)
             {me pasé, oscilo\\ \textcolor{RojoAlerta}{$\alpha$ \textbf{disminuye}}};
        \end{scope}
    \end{tikzpicture}
    \end{center}
\end{frame}

% ---------------------------------------------------------------
% 2. Pseudocódigo
% ---------------------------------------------------------------
\begin{frame}{Hypergradient Descent: Pseudocódigo}
    \scriptsize
    \begin{algorithmic}[1]
        \Require $f$, $\nabla f$, punto $\mathbf{x}$, tasa inicial $\alpha_0$,
                 \textbf{hipertasa} $\mu$
        \State $\alpha \gets \alpha_0$; \quad
               $\mathbf{g}_{\text{prev}} \gets \mathbf{0}$
               \Comment{en el paso 0 no hay gradiente anterior}
        \For{$k = 0$ \textbf{to} $k_{\max}$}
            \State $\mathbf{g} \gets \nabla f(\mathbf{x})$
            \State $\alpha \gets \alpha + \mu\,
                   \big(\mathbf{g}^{\top}\mathbf{g}_{\text{prev}}\big)$
                   \Comment{descenso sobre $\alpha$: $\alpha - \mu\,\partial f/\partial\alpha$}
            \State $\mathbf{g}_{\text{prev}} \gets \mathbf{g}$
            \State $\mathbf{x} \gets \mathbf{x} - \alpha\,\mathbf{g}$
                   \Comment{paso normal de descenso}
        \EndFor
        \State \Return $\mathbf{x}$
    \end{algorithmic}

    \vspace{0.5em}
    \begin{block}{Observaciones}
        \tiny
        \begin{itemize}
            \item El signo $+$ en la línea 4 \emph{es} un descenso: como
                  $\partial f/\partial\alpha = -\mathbf{g}^{\top}\mathbf{g}_{\text{prev}}$,
                  restarlo equivale a sumar el producto punto.
            \item Se cambia el problema de ajustar $\alpha$ por el de ajustar $\mu$,
                  pero el método es \textbf{mucho menos sensible} a $\mu$ que a $\alpha$.
            \item La misma idea se aplica sobre cualquier método de primer orden
                  (momentum, Nesterov, Adam): basta derivar respecto a su $\alpha$.
        \end{itemize}
    \end{block}
\end{frame}

% ---------------------------------------------------------------
% 3. Ejemplo paso a paso
% ---------------------------------------------------------------
\begin{frame}{Hypergradient Descent: Ejemplo Paso a Paso}
    \begin{exampleblock}{Planteamiento}
        \scriptsize
        $f(x)=x^{2}$, \ $f'(x)=2x$, \ $x^{(0)}=1$, \ $\mu=0.01$.
        Se prueban \textbf{dos tasas iniciales deliberadamente malas} para ver cómo el
        método las corrige solo.
    \end{exampleblock}

    \vspace{-0.2em}
    \begin{columns}[T]
        \begin{column}{0.49\textwidth}
            \begin{center}
            \tiny
            \textcolor{VerdeNeon}{\textbf{Caso A: $\alpha_0=0.10$ (muy chico)}}\\[0.2em]
            \begin{tabular}{c r r r}
                \toprule
                $k$ & $x^{(k)}$ & $g\cdot g_{\text{prev}}$ & $\alpha$ \\
                \midrule
                0 & $+1.0000$ & $0.000$ & $0.1000$ \\
                1 & $+0.8000$ & $3.200$ & $0.1320$ \\
                2 & $+0.5888$ & $1.884$ & $0.1508$ \\
                3 & $+0.4112$ & $0.968$ & $0.1605$ \\
                4 & $+0.2792$ & $0.459$ & $0.1651$ \\
                5 & $+0.1870$ & $0.209$ & $0.1672$ \\
                \bottomrule
            \end{tabular}
            \end{center}
            \vspace{0.2em}
            \tiny
            Los gradientes mantienen el signo $\Rightarrow$ producto $>0$
            $\Rightarrow$ $\alpha$ sube de $0.10$ a $0.167$: el método
            \textbf{acelera solo}.
        \end{column}
        \begin{column}{0.49\textwidth}
            \begin{center}
            \tiny
            \textcolor{RojoAlerta}{\textbf{Caso B: $\alpha_0=0.90$ (muy grande)}}\\[0.2em]
            \begin{tabular}{c r r r}
                \toprule
                $k$ & $x^{(k)}$ & $g\cdot g_{\text{prev}}$ & $\alpha$ \\
                \midrule
                0 & $+1.0000$ & $0.000$  & $0.9000$ \\
                1 & $-0.8000$ & $-3.200$ & $0.8680$ \\
                2 & $+0.5888$ & $-1.884$ & $0.8492$ \\
                3 & $-0.4112$ & $-0.968$ & $0.8395$ \\
                4 & $+0.2792$ & $-0.459$ & $0.8349$ \\
                5 & $-0.1870$ & $-0.209$ & $0.8328$ \\
                \bottomrule
            \end{tabular}
            \end{center}
            \vspace{0.2em}
            \tiny
            $x$ \textbf{alterna de signo} (sobrepasa el mínimo en cada paso)
            $\Rightarrow$ producto $<0$ $\Rightarrow$ $\alpha$ baja: el método
            \textbf{frena solo}.
        \end{column}
    \end{columns}

    \vspace{0.3em}
    \scriptsize
    En ambos casos $\alpha$ se mueve hacia la zona estable sin que nadie la ajuste
    a mano. Para $f=x^{2}$ el paso ideal es $\alpha=0.5$.
\end{frame}

% ---------------------------------------------------------------
% 4. Código
% ---------------------------------------------------------------
\begin{frame}[fragile]{Hypergradient Descent: Código}
\begin{lstlisting}[language=Python]
import numpy as np

def hypergradient_descent(f, grad, x, alpha=0.1, mu=1e-2,
                          n=100, eps=1e-8):
    x = np.atleast_1d(np.asarray(x, float))
    g_prev = np.zeros_like(x)        # no hay gradiente anterior
    for _ in range(n):
        g = grad(x)
        alpha = alpha + mu * (g @ g_prev)   # paso sobre alpha
        if alpha < 0:                # salvaguarda practica
            alpha = 1e-6
        g_prev = g
        x = x - alpha * g            # paso sobre x
        if np.linalg.norm(g) < eps:
            break
    return x, alpha

f    = lambda v: v @ v               # f(x) = x^2
grad = lambda v: 2*v
# hypergradient_descent(f, grad, [1.0], alpha=0.1)
#   -> x ~ 0,  alpha subio de 0.10 a ~0.17
\end{lstlisting}
\end{frame}


% B. Second-Order Methods
\section{Second-Order Methods}
\subsection{Secant Method}

\begin{frame}{Secant Method: Explicación Gráfica}
    \begin{block}{Concepto principal}
        [Explica la aproximación de la segunda derivada mediante diferencias finitas de primeras derivadas].
    \end{block}
    \begin{itemize}
        \item {[Idea gráfica 1: Recta secante a la curva de la derivada interceptando el eje cero].}
        \item {[Idea gráfica 2: Comparación geométrica frente al paso de Newton clásico].}
    \end{itemize}
\end{frame}

\begin{frame}{Secant Method: Pseudocódigo}
    \begin{algorithmic}[1]
        \State \textbf{Input:} {[Derivada \(f'\), puntos iniciales \(x^{(0)}\) y \(x^{(1)}\), tolerancia]}
        \While{[Condición de parada]}
            \State Evaluar derivadas en los dos puntos previos
            \State Aproximar curvatura secante: \(\Delta x / \Delta f'\)
            \State Calcular siguiente iteración \(x^{(k+1)}\)
        \EndWhile
        \State \textbf{Return} {[Punto estacionario]}
    \end{algorithmic}
\end{frame}

\begin{frame}{Secant Method: Ejemplo Paso a Paso}
    \begin{exampleblock}{Planteamiento del problema}
        [Define función univariable no lineal \(f(x)\) y dos valores semilla \(x_0, x_1\)].
    \end{exampleblock}
    \begin{itemize}
        \item \textbf{Paso 1:} {[Evaluación de \(f'(x_0)\) y \(f'(x_1)\)].}
        \item \textbf{Paso 2:} {[Cálculo de la secante y determinación de \(x_2\)].}
        \item \textbf{Paso 3:} {[Verificación del error respecto al óptimo].}
    \end{itemize}
\end{frame}

\begin{frame}[fragile]{Secant Method: Código}
\begin{lstlisting}[language=Python]
def secant_method(df, x0, x1, tol=1e-5, max_iter=50):
    # TODO: Bucle de la secante
    # for _ in range(max_iter):
    #     d0, d1 = df(x0), df(x1)
    #     x_next = x1 - d1 * (x1 - x0) / (d1 - d0)
    #     x0, x1 = x1, x_next
    pass
\end{lstlisting}
\end{frame}

\begin{frame}{Secant Method: Análisis de Convergencia}
    \begin{alertblock}{Tasa de Convergencia}
        [Colocar orden de convergencia superlineal (aprox. número áureo $1.618$)].
    \end{alertblock}
    \begin{itemize}
        \item {[Comparativa frente a la convergencia cuadrática de Newton].}
        \item {[Ahorro computacional al no requerir segundas derivadas analíticas].}
    \end{itemize}
\end{frame}
\subsection{Quasi-Newton Methods}

% Paleta de curvas de nivel
\colorlet{niv1}{AmarilloAviso!85!black}
\colorlet{niv2}{AmarilloAviso!50!VerdeNeon!75!black}
\colorlet{niv3}{VerdeNeon!65!black}
\colorlet{niv4}{VerdeNeon!40!CelesteBrillante!65!black}
\colorlet{niv5}{CelesteBrillante!65!black}
\colorlet{niv6}{CelesteBrillante!48!black}
\colorlet{niv7}{CelesteBrillante!34!black}
\providecommand{\reflibro}[1]{\par\vfill{\tiny\textcolor{GrisTexto!70}{\textit{Ref. libro:} #1}}}

% Comandos para las curvas de nivel de Quasi-Newton con gradiente de color
\newcommand{\qnValle}{%
    \foreach \r/\cl in {0.5/niv1, 1.0/niv2, 1.5/niv3, 2.0/niv4, 2.5/niv5, 3.0/niv6}
        \draw[\cl, thick, rotate around={35:(0,0)}] (0,0) ellipse ({\r*1.4} and {\r*0.4});%
}

\newcommand{\qnValleTrayectoria}{%
    \foreach \r/\cl in {0.3/niv1, 0.7/niv2, 1.2/niv3, 1.8/niv4, 2.5/niv5, 3.3/niv6}
        \draw[\cl, thick, rotate around={25:(0,0)}] (0,0) ellipse ({\r*1.6} and {\r*0.35});%
}

% --- FRAME 1: Introducción y Concepto ---
\begin{frame}{Quasi-Newton: Explicación Gráfica (1/3)}
    \begin{columns}[T]
        \begin{column}{0.44\textwidth}
            \centering
            \vspace{0.2cm}
            \begin{tikzpicture}[x=1.0cm, y=1.0cm, font=\scriptsize]
                \begin{scope}
                    \clip (-2.2,-2.2) rectangle (2.2,2.2);
                    
                    % Curvas de nivel con gradiente
                    \qnValle
                    
                    % Óptimo
                    \fill[white] (0,0) circle (1.5pt);
                    \draw[GrisTexto] (0,0) circle (1.5pt) node[below right, inner sep=3pt] {$x^*$};
                    
                    % Punto de inicio
                    \coordinate (X1) at (-1.5, 1.2);
                    \fill[GrisTexto!70] (X1) circle (1.6pt) node[above left, inner sep=2pt] {$x^{(k)}$};
                    
                    % Vector Gradiente Puro
                    \draw[CelesteBrillante, thick, dashed, -stealth] (X1) -- ++(0.7, 1.2) node[right, inner sep=2pt] {$-g^{(k)}$};
                    
                    % Vector Quasi-Newton
                    \draw[NaranjaBase, very thick, -stealth] (X1) -- (-0.4, 0.3) node[midway, below right, inner sep=1pt] {paso QN};
                    
                \end{scope}
                \draw[gray!50] (-2.2,-2.2) rectangle (2.2,2.2);
            \end{tikzpicture}
        \end{column}
        
        \begin{column}{0.54\textwidth}
            \scriptsize
            \begin{block}{De Newton a Quasi-Newton}
                El método de Newton exige el Hessiano exacto, calculando el paso como $x^{(k+1)} = x^{(k)} - (H^{(k)})^{-1} g^{(k)}$. En altas dimensiones, calcular e invertir esta matriz es prohibitivo.
            \end{block}
            
            \vspace{0.3cm}
            \begin{itemize}
                \setlength{\itemsep}{0.8em}
                \item \textbf{Estrategia QN:} Los métodos Quasi-Newton aproximan la inversa del Hessiano iterativamente.
                \item La regla de actualización se basa en una matriz $Q^{(k)}$ que aproxima dicha inversa en el punto $x^{(k)}$:
                \[ x^{(k+1)} \leftarrow x^{(k)} - \alpha^{(k)} Q^{(k)} g^{(k)} \]
            \end{itemize}
        \end{column}
    \end{columns}
\end{frame}

% --- FRAME 2: Descomposición del Paso (con la llave) ---
\begin{frame}{Quasi-Newton: Explicación Gráfica (2/3)}
    
    % --- FILA SUPERIOR: Ecuación y Gráfica ---
    \begin{columns}[T]
        \begin{column}{0.42\textwidth}
            \begin{block}{Iteración $k$}
                \footnotesize
                \[ \Delta x^{(k)} = -\alpha^{(k)} Q^{(k)} g^{(k)} \]
            \end{block}
        \end{column}
        
        \begin{column}{0.56\textwidth}
            \centering
            \begin{tikzpicture}[x=0.95cm, y=0.95cm, font=\scriptsize]
                \begin{scope}
                    \clip (-2.2,-2.2) rectangle (2.2,2.2);
                    
                    % Curvas de nivel con gradiente
                    \qnValle
                    
                    \coordinate (Xk) at (-1.5, 1.2);
                    
                    % Descomposición de vectores
                    \draw[CelesteBrillante, very thick, -stealth] (Xk) -- (-0.8, 2.4) node[midway, above left, inner sep=2pt] {$-g^{(k)}$};
                    \draw[VerdeNeon, very thick, dotted, -stealth] (-0.8, 2.4) -- (-0.4, 0.3) node[midway, right, inner sep=2pt] {corrección $Q^{(k)}$};
                    \draw[NaranjaBase, ultra thick, -stealth] (Xk) -- (-0.4, 0.3);
                    
                    \fill[GrisTexto!70] (Xk) circle (1.6pt) node[left, inner sep=3pt] {$x^{(k)}$};
                    \fill[NaranjaBase] (-0.4, 0.3) circle (1.8pt) node[right, inner sep=3pt, text=GrisTexto] {$x^{(k+1)}$};
                \end{scope}
                \draw[gray!50] (-2.2,-2.2) rectangle (2.2,2.2);
            \end{tikzpicture}
        \end{column}
    \end{columns}
    
    \vspace{-0.2cm}
    
    % --- FILA INFERIOR: Textos explicativos alineados ---
    \begin{columns}[T]
        \begin{column}{0.4\textwidth}
            \scriptsize
            \begin{itemize}
                \setlength{\itemsep}{0.5em}
                \item \tikz[remember picture, overlay] \coordinate (item1top); \textcolor{CelesteBrillante}{$-g^{(k)}$ (Gradiente)}: Apunta a la máxima pendiente local, provocando rebote.
                \item \textcolor{VerdeNeon}{$Q^{(k)}$ (Aprox. Hessiano)}: Multiplica al gradiente y al factor escalar $\alpha^{(k)}$ para aproximar el paso de Newton.\tikz[remember picture, overlay] \coordinate (item2bot);
            \end{itemize}
        \end{column}
        
        \begin{column}{0.6\textwidth}
            \scriptsize
            \vspace{0.6cm}
            \begin{itemize}
                \item \tikz[remember picture, overlay] \coordinate (item3); \textcolor{NaranjaBase}{Resultado}: El paso resultante se reorienta de forma eficiente hacia el óptimo.
            \end{itemize}
        \end{column}
    \end{columns}
    
    % --- Llave conectando las columnas inferiores ---
    \begin{tikzpicture}[remember picture, overlay]
        \coordinate (brace_x) at ([xshift=-26pt]item3);
        \draw[decorate, decoration={brace, amplitude=6pt}, thick, GrisTexto] 
            ([yshift=2ex]brace_x |- item1top) -- ([yshift=-0.5ex]brace_x |- item2bot);
    \end{tikzpicture}
    
\end{frame}

% --- FRAME 3: Trayectorias y Convergencia ---
\begin{frame}{Quasi-Newton: Explicación Gráfica (3/3)}
    \begin{columns}[T]
        \begin{column}{0.38\textwidth}
            \centering
            % Gráfico de Trayectorias
            \begin{tikzpicture}[x=0.9cm, y=0.9cm, font=\scriptsize]
                \begin{scope}
                    \clip (-2.5,-2.2) rectangle (2.5,1.5);
                    
                    % Curvas de nivel con gradiente para la trayectoria
                    \qnValleTrayectoria
                    
                    \fill[white] (0,0) circle (1.5pt);
                    \draw[GrisTexto] (0,0) circle (1.5pt) node[right, inner sep=3pt] {$x^*$};
                    
                    % Trayectoria GD (Zig-Zag)
                    \draw[CelesteBrillante, thick] (-2.2, 1.0) -- (-1.8, -0.6) -- (-1.0, 0.4) -- (-0.5, -0.3) -- (-0.1, 0.2) -- (0,0);
                    \fill[CelesteBrillante] (-2.2, 1.0) circle (1.2pt) node[above right, text=GrisTexto] {$x^{(1)}$};
                    \foreach \p in {(-1.8, -0.6), (-1.0, 0.4), (-0.5, -0.3), (-0.1, 0.2)} { \fill[CelesteBrillante] \p circle (1pt); }
                    
                    % Trayectoria Quasi-Newton
                    \draw[NaranjaBase, very thick] (-2.2, 1.0) -- (-0.8, 0.2) -- (-0.1, 0.05) -- (0,0);
                    \foreach \p in {(-0.8, 0.2), (-0.1, 0.05)} { \fill[NaranjaBase] \p circle (1.5pt); }
                \end{scope}
                \draw[gray!50] (-2.5,-2.2) rectangle (2.5,1.5);
                
                % Leyenda integrada
                \node[text=NaranjaBase, right] at (-2.4, -2.5) {\textbf{---} Quasi-Newton};
                \node[text=CelesteBrillante, right] at (0.2, -2.5) {\textbf{---} GD};
            \end{tikzpicture}
            
            \vspace{0.2cm}
            % Gráfico de Barras Corregido
            \begin{tikzpicture}[font=\scriptsize]
                \node[text=GrisTexto, font=\footnotesize] at (2.2, 1.8) {Distancia al mínimo ($||x_k - x^*||$)};
                \draw[gray] (0,0) -- (4.5,0);
                
                % Iteraciones 1 a 6
                \foreach \x/\qn/\gd in {1/1.5/1.5, 2/0.7/1.3, 3/0.15/1.1, 4/0.02/0.9, 5/0/0.75, 6/0/0.6} {
                    \fill[CelesteBrillante] ({\x*0.65+0.15},0) rectangle ++(0.2,\gd);
                    \fill[NaranjaBase] ({\x*0.65-0.1},0) rectangle ++(0.2,\qn);
                    \node[below] at ({\x*0.65+0.12},0) {\x};
                }
                \node[below] at (2.2, -0.3) {iteración $k$};
            \end{tikzpicture}
        \end{column}

        \scriptsize
        \begin{column}{0.60\textwidth}
            \begin{block}{Actualización BFGS}
                Se definen los cambios $\gamma^{(k+1)} \equiv g^{(k+1)} - g^{(k)}$ y $\delta^{(k+1)} \equiv x^{(k+1)} - x^{(k)}$. La aproximación se enriquece iterativamente:
                \vspace{-0.1cm}
                \[ Q \leftarrow Q - \left(\frac{\delta\gamma^\top Q + Q\gamma\delta^\top}{\delta^\top\gamma}\right) + \left(1 + \frac{\gamma^\top Q\gamma}{\delta^\top\gamma}\right)\frac{\delta\delta^\top}{\delta^\top\gamma} \]
            \end{block}
            
            \vspace{0.2cm}
            \begin{itemize}
                \setlength{\itemsep}{0.8em}
                \item \textbf{Costo de Memoria:} Aunque acelera la convergencia drásticamente respecto al GD, la matriz densa $Q$ demanda grandes cantidades de memoria espacial.
                \item \textbf{L-BFGS:} El método L-BFGS fue diseñado para aproximar a BFGS reduciendo este impacto. La matriz inversa completa se sustituye almacenando únicamente los últimos $m$ valores de $\delta$ y $\gamma$.
            \end{itemize}
        \end{column}
    \end{columns}
\end{frame}

% ---------------------------------------------------------------
% 2. Pseudocódigo
% ---------------------------------------------------------------
\begin{frame}{Quasi-Newton (BFGS): Pseudocódigo}
    \scriptsize
    \begin{algorithmic}[1]
        \Require $f$, $\nabla f$, $x^{(1)}$, tolerancia $\epsilon$, $k_{\max}$
        \State $Q \gets I$ \Comment{Matriz identidad: aprox. inicial de la inversa de la Hessiana}
        \For{$k = 1, \dots, k_{\max}$}
            \State $g \gets \nabla f(x)$
            \State \textbf{if} $\lVert g \rVert < \epsilon$ \textbf{then break}
            \State $d \gets -Q g$ \Comment{Dirección de descenso casi-Newton (6.22)}
            \State $x_{\text{new}} \gets \text{line\_search}(f, x, d)$ \Comment{Búsqueda de línea a lo largo de $d$}
            \State $g_{\text{new}} \gets \nabla f(x_{\text{new}})$
            \State $\delta \gets x_{\text{new}} - x$ \Comment{(6.24)}
            \State $\gamma \gets g_{\text{new}} - g$ \Comment{(6.23)}
            \State \textbf{if} $\delta^\top \gamma \approx 0$ \textbf{then break} \Comment{Prevenir división por cero}
            \State $Q \gets Q - \frac{\delta\gamma^\top Q + Q\gamma\delta^\top}{\delta^\top\gamma} + \left(1 + \frac{\gamma^\top Q\gamma}{\delta^\top\gamma}\right)\frac{\delta\delta^\top}{\delta^\top\gamma}$ \Comment{Actualización BFGS (6.26)}
            \State $x \gets x_{\text{new}}$
        \EndFor
        \State \Return $x$
    \end{algorithmic}

    \vspace{1em}
    \scriptsize
    \begin{itemize}
        \item Aproxima la inversa de la Hessiana ($H^{-1}$) usando solo información de gradientes pasados.
        \item Mantiene $Q$ simétrica y definida positiva en cada paso.
        \item Existen otras variantes como DFP (Ec. 6.25), pero BFGS es el estándar más robusto.
    \end{itemize}
    \reflibro{Algoritmo~6.6, Ec.~(6.22)--(6.26).}
\end{frame}

% ---------------------------------------------------------------
% 3. Ejemplo paso a paso
% ---------------------------------------------------------------
\begin{frame}{Quasi-Newton (BFGS): Ejemplo Paso a Paso (1/3)}
    \begin{exampleblock}{Planteamiento}
        \scriptsize
        $f(x)=\tfrac12x_1^2+5x_2^2$, \quad $\nabla f=(x_1,\ 10x_2)$, \quad
        $x^{(1)}=(-4,\,1)$, \quad $\alpha=0.15$
    \end{exampleblock}

    \begin{columns}[T]
        \begin{column}{0.5\textwidth}
            \begin{block}{Inicialización}
                \scriptsize
                $Q^{(1)}=I=\begin{pmatrix} 1 & 0 \\ 0 & 1 \end{pmatrix}$
            \end{block}
            \begin{block}{Iteración $k=1$}
                \scriptsize
                $g^{(1)}=(-4,\ 10)$\\[2pt]
                $d^{(1)}=-Q^{(1)}g^{(1)}=(4,\ -10)$\\[2pt]
                $x^{(2)}=x^{(1)} + 0.15\, d^{(1)} = (-3.40,\ -0.50)$\\[2pt]
                $\mathbf{f(x^{(2)})=7.030}$
            \end{block}
            {\scriptsize \textbf{Nota:} Al ser $Q^{(1)}=I$, el primer paso es idéntico a usar Descenso de Gradiente (GD).\par}
        \end{column}
        \begin{column}{0.48\textwidth}
            \centering
            \begin{tikzpicture}[x=1.05cm, y=1.05cm, font=\scriptsize]
                \begin{scope}
                    \clip (-4.4,-1.3) rectangle (0.6,1.35);
                    \momElipses
                    \draw[NaranjaBase, very thick, -stealth] (-4,1) -- (-3.40,-0.50);
                \end{scope}
                \draw[gray!50] (-4.4,-1.3) rectangle (0.6,1.35);
                \fill[GrisTexto] (-4,1) circle (1.8pt) node[right, inner sep=3pt] {$x^{(1)}$};
                \fill[GrisTexto] (-3.40,-0.50) circle (1.8pt) node[right, inner sep=3pt] {$x^{(2)}$};
                \fill[white] (0,0) circle (1.6pt) node[above, inner sep=2pt] {$x^*$};
            \end{tikzpicture}
        \end{column}
    \end{columns}
\end{frame}

\begin{frame}{Quasi-Newton (BFGS): Ejemplo Paso a Paso (2/3)}
    \begin{columns}[T]
        \begin{column}{0.56\textwidth}
            \begin{block}{Iteración $k=2$}
                \scriptsize
                $g^{(2)}=(-3.40,\ -5.00)$\\[2pt]
                $\delta^{(2)}=x^{(2)}-x^{(1)}=(0.60,\ -1.50)$\\[2pt]
                $\gamma^{(2)}=g^{(2)}-g^{(1)}=(0.60,\ -15.00)$\\[2pt]
                Actualizamos $Q^{(2)}$ usando la Ec. (6.26)\\[2pt]
                $d^{(2)}=-Q^{(2)}g^{(2)} \approx (3.90,\ 0.52)$\\[2pt]
                $x^{(3)}=x^{(2)}+0.15\, d^{(2)} = (-2.81,\ -0.42)$\\[2pt]
                $\mathbf{f(x^{(3)})=4.851}$
            \end{block}
            \small
            \begin{itemize}
                \item $Q$ se moldea capturando la curvatura.
                \item La dirección de descenso se desvía del gradiente negativo y apunta de forma más directa hacia el mínimo.
            \end{itemize}
        \end{column}
        \begin{column}{0.42\textwidth}
            \centering
            \begin{tikzpicture}[x=1.45cm, y=1.45cm, font=\scriptsize]
                \begin{scope}
                    \clip (-4.25,-1.75) rectangle (-1.9,1.15);
                    \momElipses
                    \draw[GrisTexto!60, thick, -stealth] (-4,1) -- (-3.40,-0.50);
                    \draw[NaranjaBase, very thick, -stealth] (-3.40,-0.50) -- (-2.81,-0.42);
                \end{scope}
                \draw[gray!50] (-4.25,-1.75) rectangle (-1.9,1.15);
                \fill[GrisTexto] (-4,1) circle (1.6pt) node[right, inner sep=3pt] {$x^{(1)}$};
                \fill[GrisTexto] (-3.40,-0.50) circle (1.6pt) node[left, inner sep=2pt] {$x^{(2)}$};
                \fill[VerdeNeon] (-2.81,-0.42) circle (1.8pt) node[right, inner sep=2pt, text=GrisTexto] {$x^{(3)}$};
            \end{tikzpicture}
        \end{column}
    \end{columns}
\end{frame}

\begin{frame}{Quasi-Newton (BFGS): Ejemplo Paso a Paso (3/3)}
    \begin{columns}[T]
        \begin{column}{0.56\textwidth}
            \begin{block}{Iteración $k=3$}
                \scriptsize
                $g^{(3)}=(-2.81,\ -4.22)$\\[3pt]
                Actualizamos $Q^{(3)}$ \dots\\[3pt]
                $x^{(4)}=(-2.39,\ -0.36)$\\[3pt]
                $\mathbf{f(x^{(4)})=3.505}$
            \end{block}
            \begin{exampleblock}{Resultado final ($k_{\max}=60$)}
                \scriptsize
                $x^{(61)} = (-0.0002,\ 0.0000)$\\[3pt]
                $f(x^{(61)}) = 0.0000$
            \end{exampleblock}
        \end{column}
        \begin{column}{0.42\textwidth}

            \centering
            \begin{tikzpicture}[x=0.95cm, y=0.95cm, font=\scriptsize]
                \begin{scope}
                    \clip (-4.4,-1.15) rectangle (0.5,1.3);
                    \momElipses
                    \draw[NaranjaBase, very thick] plot coordinates {(-4.00,1.00) (-3.40,-0.50) (-2.81,-0.42) (-2.39,-0.36) (-2.03,-0.30) (-1.73,-0.26) (-1.47,-0.22)};
                    \foreach \p in {(-3.40,-0.50), (-2.81,-0.42), (-2.39,-0.36), (-2.03,-0.30), (-1.73,-0.26), (-1.47,-0.22)}
                        \fill[NaranjaBase] \p circle (1.4pt);
                \end{scope}
                \draw[gray!50] (-4.4,-1.15) rectangle (0.5,1.3);
                \fill[white] (-4,1) circle (1.6pt) node[right, inner sep=3pt, text=GrisTexto] {$x^{(1)}$};
                \fill[white] (0,0) circle (1.4pt);
                \draw[NaranjaBase, very thick] (-4.3,-1.45) -- (-3.9,-1.45) node[right, text=GrisTexto] {BFGS};
            \end{tikzpicture}

            {\tiny Primeras 7 iteraciones de BFGS.\par}
        \end{column}
    \end{columns}

    \vspace{0.3em}
    \centering
    {\scriptsize
    \begin{tabular}{l c c c c c c c}
        \toprule
        $k$ & 0 & 1 & 2 & 3 & 4 & 5 & 6 \\
        \midrule
        $f(x^{(k+1)})$ BFGS & 13.00 & 7.030 & 4.851 & 3.505 & 2.532 & 1.830 & 1.322 \\
        \bottomrule
    \end{tabular}\par}
\end{frame}

% ---------------------------------------------------------------
% 4. Código
% ---------------------------------------------------------------
\begin{frame}[fragile]{Quasi-Newton (BFGS): Código}
\begin{lstlisting}[language=Python]
import numpy as np

def bfgs(f, grad, x, alpha=0.1, k_max=100, eps_tol=1e-6):
    m = len(x)
    Q = np.eye(m)  # Inicializacion de Q^{(1)} = I
    for k in range(1, k_max + 1):
        g = grad(x)
        if np.linalg.norm(g) < eps_tol:
            break         
        d = -Q @ g                        # (6.22) Direccion
        x_new = x + alpha * d             # Paso (reemplaza line_search)
        g_new = grad(x_new)
        delta = x_new - x                 # (6.24)
        gamma = g_new - g                 # (6.23)
        denom = np.dot(delta, gamma)
        if denom == 0:
            break    
        # Actualizacion BFGS (6.26)
        term1 = (np.outer(delta, gamma) @ Q + Q @ np.outer(gamma, delta)) / denom
        term2 = (1 + np.dot(gamma, Q @ gamma) / denom) * (np.outer(delta, delta) / denom)
        Q = Q - term1 + term2
        x = x_new
    return x, k

# Ejemplo: f(x) = 0.5*x1^2 + 5*x2^2   ->   grad f = (x1, 10*x2)
f = lambda x: 0.5*x[0]**2 + 5*x[1]**2
grad = lambda x: np.array([x[0], 10*x[1]])
x0 = np.array([-4.0, 1.0])
bfgs(f, grad, x0, alpha=0.1, k_max=100)
\end{lstlisting}
\end{frame}

% Bibliografía
\section{Bibliografía}

\begin{frame}{Bibliografía}
    \begin{thebibliography}{99}
        \bibitem{kochenderfer2026}
        Kochenderfer, M. J., \& Wheeler, T. A. (2026). \\
        \textit{Algorithms for optimization} (2nd ed.). \\
        The MIT Press.
    \end{thebibliography}
\end{frame}

\end{document}